% !TEX program = xelatex
% Статья для ПЗ-3. Описаны метод и план проверки; эксперимент не проведён.
\documentclass{misisarticle}
\usepackage{balance}
\urlstyle{same}

\grnti{28.23}
\asjc{1702, 1706}
\oecd{1.2}
\title{Методы формирования персонализированных объяснений в RecSys с использованием LLM и структурированных моделей пользователей}
% Контактный email не указан: не подставлять адрес из шаблона.
\misisauthor{В.~С.~Сердешнов}{кафедра инженерной кибернетики НИТУ «МИСиС»}{Москва, Россия}{}

\setabstract{Рассматривается задача формирования текстовых объяснений результатов рекомендательной системы. Предлагается двухэтапная схема: сначала из структурированного профиля пользователя и метаданных объекта отбираются проверяемые основания, затем большая языковая модель формирует по ним текст. Для каждого основания сохраняются исходные записи и правило его получения. Сформулированы математическая постановка отбора оснований и показатели подтверждаемости и персонализированности объяснений. Разработан протокол сравнения шаблонного и признакового методов, языковой модели без профиля и модели со структурированным профилем при фиксированных рекомендациях. Основным источником данных выбран MTS Kion; для пилота предусмотрена ручная разметка объяснений. Работа представляет описание метода и проект экспериментальной проверки; количественные результаты и выводы о превосходстве подхода на данном этапе отсутствуют.}
\setkeywords{рекомендательные системы, персонализированные объяснения, большие языковые модели, структурированные профили пользователей, проверяемые основания, фактическая достоверность}
\hypersetup{pdftitle={Персонализированные объяснения рекомендаций},pdfauthor={В. С. Сердешнов}}

\begin{document}
\maketitle

\section{Введение}

Рекомендательная система выбирает объекты из каталога на основе информации о пользователе и его взаимодействиях с сервисом. Дополнительный текст может пояснить, какие свойства объекта связаны с доступной историей пользователя. При этом цели объяснения различаются: прозрачность работы системы, помощь в выборе, убедительность и доверие не являются одной и той же характеристикой~\cite{survey}. Поэтому качество текста нельзя свести к его связности или положительной реакции читателя.

Большие языковые модели (LLM) позволяют формировать разнообразные объяснения, но содержательная правдоподобность не гарантирует подтверждаемости утверждений. Например, из наблюдения «пользователь просмотрел несколько драм» не следует утверждение «пользователь предпочитает драмы всем другим жанрам». В первом случае описываются данные, во втором добавляется интерпретация, для которой может не быть оснований. Такая подмена особенно существенна при работе с неявной обратной связью: просмотр не равнозначен явной положительной оценке.

В работе исследуется формирование объяснений уже полученных рекомендаций. Предлагается использовать структурированный профиль с происхождением признаков и ограниченный набор проверяемых оснований. Исследовательский вопрос состоит в том, как такая организация входных данных влияет на подтверждаемость и персонализированность текста относительно простых методов и генерации без профиля.

Методическая часть работы включает представление оснований, постановку их отбора и протокол оценки на одинаковых рекомендациях. Предлагаемые элементы не объявляются доказанно новыми или эффективными до сравнительного исследования. Настоящая статья описывает постановку, методы и план проверки, создавая основу для последующего эксперимента.

\section{Анализ связанных работ}

В обзоре Zhang и Chen~\cite{survey} рассматриваются различные способы объяснения рекомендаций, включая методы, встроенные в модель, и объяснения, формируемые после ранжирования. Для рассматриваемой задачи существенно различать описание наблюдаемых свойств пары «пользователь~--- объект» и раскрытие механизма принятия решения рекомендателем.

Li, Zhang и Chen~\cite{pepler} предлагают персонализированное обучение промптов для генерации объяснений. Авторы исследуют дискретные представления идентификаторов и непрерывные векторные промпты, позволяющие связать пользователя и объект с предобученной языковой моделью. Для нашей постановки важна возможность персонализировать текст; отдельно требуется проверить, какие его утверждения подтверждаются доступными записями.

Ramos и соавторы~\cite{upr} используют профили на естественном языке, полученные из отзывов, для построения рекомендаций и исследуют изменение выдачи при редактировании профиля. Такой профиль доступен для просмотра и корректировки человеком. В настоящей работе предлагается хранить характеристики в явной структуре вместе со свидетельствами и использовать их для объяснения фиксированной выдачи. Это другая экспериментальная постановка: не замена рекомендателя текстовым профилем, а проверка качества текста по разрешённым фактам.

В XRec~\cite{xrec} коллаборативные сигналы связываются с LLM посредством адаптера. Подход позволяет включать информацию о взаимодействиях в генерацию. Здесь рассматривается вариант без обучения коллаборативного адаптера: модель получает явные признаки профиля и основания, которые можно проверить по исходным данным. Поэтому XRec и PEPLER выступают прежде всего содержательными ориентирами обзора; их полноценное воспроизведение не входит в первоначальное сравнение четырёх способов.

В сборнике кафедры работа Соседки и Ибрагимова~\cite{sosedka} посвящена анализу аудитории Telegram в контексте рекомендаций с использованием LLM, а работа Береснева~\cite{beresnev}~--- предобученным моделям для корпоративных чат-ботов. Они рассматривают смежные прикладные задачи. Публикация Марчуковой и Кожаринова~\cite{kozharinov} относится к прогнозированию и принятию решений; её тематическая связь косвенная, и она не используется как подтверждение эффективности LLM-объяснений.

В предлагаемой схеме предметом проверки является связь между структурой профиля, происхождением оснований и свойствами итогового текста. Наличие соответствующей записи или корректной ссылки на неё само по себе не гарантирует верности формулировки. Эта связь должна оцениваться содержательной разметкой.

\section{Постановка задачи}

\subsection{Входные данные и граница исследования}

Пусть $U$~--- множество пользователей, $I$~--- каталог объектов, $t_0$~--- момент отсечения истории. Доступные взаимодействия задаются множеством
\begin{equation}
 D_{<t_0}=\{(u,i,t,w):t<t_0\},
 \label{eq:history}
\end{equation}
где $w$~--- характеристика просмотра. Для объекта $i$ известны метаданные $M_i$, для пользователя $u$~--- история $H_u\subset D_{<t_0}$. Отдельный рекомендатель по этим данным формирует список $R_K(u)$ из $K$ ещё не просмотренных объектов.

По истории вычисляется профиль $P_u=f(H_u,M)$, где $M=\{M_i:i\in I\}$. Из профиля и метаданных рекомендуемого объекта строится множество кандидатов в основания $B_{ui}$. Каждое основание представляет собой запись
\begin{equation}
 b=(\mathrm{id},\,\mathrm{predicate},\,\mathrm{value},\,S_b),
 \label{eq:evidence}
\end{equation}
где $S_b$ содержит ссылки на исходные записи и правило вычисления. Например, основание может указывать число просмотренных объектов определённого жанра и список этих объектов. Метаданные рекомендуемого произведения и вычисленные характеристики истории хранятся раздельно.

Требуется выбрать $E_{ui}\subseteq B_{ui}$ и сформировать объяснение $y_{ui}$ для каждой оцениваемой пары $(u,i)$, $i\in R_K(u)$. Ранжирование, отбор оснований и генерация являются связанными подзадачами, но оцениваются отдельно. Корректное описание сходства объекта с историей не доказывает, что именно это сходство определило внутренний балл рекомендателя. Далее под подтверждаемостью понимается опора текста на данные, а не причинная верность объяснения механизму ранжирования.

\subsection{Отбор оснований и условная генерация}

Для каждого основания задаются неотрицательная оценка релевантности $s(b,u,i)$ и размер $\ell(b)$ в токенах сериализованного представления. Допустимыми считаются основания, проверенные по истории до $t_0$ и метаданным. Их отбор формулируется как задача
\begin{equation}
 \begin{aligned}
 E^*_{ui}&\in\mathop{\mathrm{arg\,max}}_{E\subseteq B_{ui}}
                 \sum_{b\in E}s(b,u,i),\\
 |E|&\leq k,\qquad \sum_{b\in E}\ell(b)\leq L_E .
 \end{aligned}
 \label{eq:selection}
\end{equation}
Здесь $k$ ограничивает число оснований, $L_E$~--- объём контекста. Формула задаёт критерий отбора, но не требует сложного оптимизатора: для небольшого числа жанровых оснований можно перебрать допустимые подмножества. Равные оценки разрешаются фиксированным порядком идентификаторов.

LLM с параметрами $\theta$ получает выбранные основания и инструкцию $C$:
\begin{equation}
 y_{ui}\sim p_\theta(\,\cdot\mid M_i,P_u,E^*_{ui},C).
 \label{eq:generator}
\end{equation}
Пусть $A(y)$~--- множество атомарных проверяемых утверждений текста, а $z(a,E)$ равен единице, если набор $E$ поддерживает утверждение $a$. Целевое условие имеет вид
\begin{equation}
 \forall a\in A(y_{ui}):\quad z(a,E^*_{ui})=1.
 \label{eq:grounding}
\end{equation}
Оно выражает требование к результату, а не математическую гарантию работы генератора. Проверка JSON-схемы, наличие идентификатора основания и инструкция в промпте не заменяют оценку смысловой поддержки. В пилоте $z$ определяется ручной разметкой.

\section{Описание методов}

\subsection{Данные и формирование профиля}

Основным источником выбран открытый MTS Kion~\cite{kion}. Он содержит неявные взаимодействия с параметрами просмотра, характеристики пользователей и метаданные контента, включая жанры и описания. Эти сведения позволяют строить профиль и проверять текст без добавления внешних знаний LLM. Тестовые взаимодействия не используются ни при вычислении профиля, ни при выборе оснований.

В первоначальном варианте персонализация опирается на историю просмотра. Демографические признаки не включаются в текст автоматически: возрастная группа или наличие детей не доказывают предпочтение конкретного жанра. Неполный просмотр также не трактуется как явная отрицательная оценка.

Пусть $H^+_u$ содержит уникальные объекты с положительной зарегистрированной долей просмотра и известным непустым набором жанров. Для $|H^+_u|>0$ определим распределение
\begin{equation}
 \pi_u(g)=\frac{1}{|H^+_u|}
 \sum_{j\in H^+_u}\frac{\mathbf{1}[g\in G_j]}{|G_j|},
 \label{eq:profile}
\end{equation}
где $G_j$~--- непустое множество жанров объекта $j$. История содержит уникальные объекты, а записи без жанра исключаются из данного агрегата. Дробный вклад многожанрового объекта обеспечивает $\sum_g\pi_u(g)=1$. Это распределение наблюдавшихся жанров, а не вероятность истинного предпочтения пользователя.

Профиль включает $\pi_u$, число объектов истории и их идентификаторы. Для жанра $g\in G_i$ основание содержит $\pi_u(g)$, количество наблюдений и ссылки на просмотренные объекты с этим жанром. В первоначальном отборе можно принять $s(b,u,i)=\pi_u(g)$ для такого основания. Дополнительные основания по актёрам и режиссёрам рассматриваются как расширение, а не обязательная часть пилота. При отсутствии пригодной истории метод возвращает неперсонализированное описание с явным указанием недостатка данных.

Для следующего этапа планируется перенос на одну категорию Amazon Reviews 2023~\cite{amazon}. Второй домен позволит проверить применимость схемы вне видеокаталога. Отзывы при этом рассматриваются как источник наблюдений и возможных аспектов, а не как безусловно истинные объяснения решений рекомендателя. MovieLens может использоваться для дополнительного сравнения внутри домена кино.

\subsection{Фиксированный рекомендатель}

В качестве базовой модели предлагается использовать матричную факторизацию для неявной обратной связи~\cite{als}. Для наблюдаемых объектов $r_{ui}$ задаётся нормированной долей просмотра, для ненаблюдаемых~--- нулём. При подготовке данных доля приводится к диапазону $[0,1]$; пропуски не подменяются отрицательными оценками. Для повторных записей одного объекта до $t_0$ используется максимальная корректная доля просмотра. Определим
\begin{equation}
 p_{ui}=\mathbf{1}[r_{ui}>0],\qquad c_{ui}=1+\alpha r_{ui},
 \label{eq:implicit}
\end{equation}
где $\alpha>0$. Нулевое значение $p_{ui}$ означает отсутствие зарегистрированного положительного сигнала, а не доказанную неприязнь к объекту. Выбор доли просмотра для $r_{ui}$ является адаптацией схемы к Kion и требует проверки чувствительности.

Пользовательские векторы $x_u\in\mathbb{R}^d$ и векторы объектов $v_i\in\mathbb{R}^d$ определяются минимизацией
\begin{equation}
 \begin{aligned}
 \mathcal{L}(X,V)={}&\sum_{u\in U}\sum_{i\in I}
 c_{ui}(p_{ui}-x_u^{\mathsf T}v_i)^2\\
 &+\lambda\left(\sum_{u\in U}\|x_u\|_2^2
                  +\sum_{i\in I}\|v_i\|_2^2\right),
 \end{aligned}
 \label{eq:als}
\end{equation}
где $d$~--- размерность, $\lambda>0$~--- коэффициент регуляризации. Метод ALS поочерёдно оптимизирует векторы пользователей и объектов при фиксированной другой группе. Гиперпараметры выбираются на валидационной части по качеству ранжирования, после чего рекомендации сохраняются и одинаково используются всеми методами объяснения. Популярностный рекомендатель служит дополнительным контролем качества выдачи. Сам ALS не рассматривается как предлагаемый научный вклад.

\subsection{Четыре способа объяснения}

Для каждой пары используются следующие варианты.

\textit{T: шаблонный.} Фиксированная конструкция заполняется свойствами рекомендуемого объекта, например его жанрами. Метод не получает профиль и служит контролем подтверждаемости неперсонализированного текста.

\textit{F: признаковый.} Из $E^*_{ui}$ выбираются основания, связывающие объект с историей. Детерминированный шаблон сообщает наблюдение, например: «В вашей истории есть просмотренные произведения этого жанра». Численные формулировки допускаются только при наличии соответствующего агрегата. Утверждения о мотивах или эмоциях пользователя не добавляются.

\textit{L0: LLM без профиля.} Модель получает $M_i$ и инструкцию, но не историю и не пользовательские признаки. Этот вариант проверяет генерацию описания объекта без персонального контекста.

\textit{LP: LLM с профилем.} Та же модель получает $P_u$, $M_i$ и $E^*_{ui}$ и формирует объяснение с идентификаторами использованных оснований. Ограничение на длину текста и параметры генерации фиксируются для обоих LLM-вариантов. Профиль содержит только разрешённые агрегаты; подробная история, не попавшая в набор свидетельств, в промпт не передаётся.

Сравнение LP и L0 оценивает совместное добавление профиля и персональных оснований. Оно не позволяет приписать эффект одной JSON-структуре. Для проверки именно представления предусмотрена дополнительная абляция: передать той же модели те же факты свободным текстом с сохранением их содержания. Для сравнения языковой генерации и детерминированной вербализации методы LP и F используют одинаковые основания.

\subsection{Контроль утверждений}

Генератор возвращает текст и ссылки на использованные основания в структурированном ответе. Техническая проверка контролирует формат и существование идентификаторов. Семантическая проверка устанавливает, соответствует ли утверждение источнику: правильная ссылка может сопровождать преувеличение или неверное обобщение.

В пилоте оцениваются исходные ответы до смысловой фильтрации, включая технически некорректные и пустые ответы в общих показателях покрытия. Это исключает скрытое улучшение результата за счёт удаления неудачных примеров. Если затем вводится фильтр, его эффект измеряется отдельно до и после применения. Демонстрационные фразы в описании методов не являются экспериментальными наблюдениями.

\section{Протокол экспериментальной проверки}

\subsection{Разбиение и пилотная выборка}

Используется хронологическое разбиение на обучение, валидацию и тест. Правила подготовки данных и значения гиперпараметров фиксируются до оценки на тесте. Качество ранжирования оценивается отдельно, например по NDCG@10 и Recall@10. Для объяснений после настройки рекомендателя сохраняется один и тот же список пар и входных данных для всех вариантов.

В пилоте планируется выбрать 30--50 пользователей с пригодной для построения профиля историей и по одному объекту из их топ-10. Четыре метода дадут 120--200 текстов. Одна пара на пользователя уменьшает зависимость наблюдений в первоначальном сравнении. Такая выборка предназначена для отладки разметки и оценки неопределённости, а не заранее гарантирует статистическую мощность. Правила отбора пользователей и объектов, число наблюдений и причины исключения публикуются вместе с результатами.

\subsection{Разметка и показатели}

Для всех методов разрешённым источником проверки служит один и тот же набор: история пользователя до $t_0$ и метаданные объекта. Он фиксируется независимо от того, какую часть сведений получает генератор. Текст делится на минимальные проверяемые утверждения. Для каждого $a\in A(y)$ размечаются два бинарных признака: $z_a$~--- поддержка разрешёнными данными, $h_a$~--- корректная связь утверждения с профилем пользователя. По определению $h_a\leq z_a$. Субъективные характеристики, например «увлекательный», учитываются отдельно и не повышают число подтверждённых фактов.

Для $N(y)=|A(y)|>0$ определяются
\begin{equation}
 F(y)=\frac{\sum_{a\in A(y)}z_a}{N(y)},\qquad
 P(y)=\frac{\sum_{a\in A(y)}h_a}{N(y)}.
 \label{eq:metrics}
\end{equation}
Показатель $F$ характеризует подтверждаемость, $P$~--- долю корректных персонализированных утверждений. При $N(y)=0$ обе доли не определяются. Пустое объяснение не получает максимальной оценки. Покрытие фактическими объяснениями метода $m$ задаётся как
\begin{equation}
 C_m=\frac{\sum_{j=1}^{n}\mathbf{1}[N(y_{jm})>0]}{n},
 \label{eq:coverage}
\end{equation}
где $n$~--- число всех исходно выбранных пар, включая случаи отказа и технической ошибки. Дополнительно сообщаются доля текстов с корректной персональной связью, число утверждений, длина ответа и задержка генерации. Без этих величин высокий $F$ может отражать только короткие общие формулировки.

Для каждого метода публикуется среднее по объяснениям с определёнными долями и их число. Парные различия вычисляются на пересечении пар, для которых показатель определён у обоих методов; размер пересечения сообщается явно. Покрытие показывается по всей исходной выборке, чтобы отказы не исчезали из оценки.

Два разметчика независимо оценивают тексты, не видя названия метода, но имея доступ к разрешённым основаниям. Для сопоставимой разметки атомарные утверждения выделяются по единой инструкции до оценки поддержки. Согласованность бинарных оценок измеряется каппой Коэна~\cite{cohen}; спорные случаи затем обсуждаются. LLM-судья может применяться как вспомогательный инструмент только после сопоставления с человеческими оценками.

\subsection{Сравнение и воспроизводимость}

Основное сравнение~--- LP против L0 по $F$; $P$ и покрытие рассматриваются совместно с ним. Направленная исследовательская гипотеза состоит в увеличении подтверждаемости при добавлении профиля и оснований. Для парных разностей
\begin{equation}
 d_j=F(y_{j,\mathrm{LP}})-F(y_{j,\mathrm{L0}})
 \label{eq:paired}
\end{equation}
предлагается сообщать среднее различие и доверительный интервал с повторной выборкой пользователей. Критерий Вилкоксона~\cite{wilcoxon} может использоваться дополнительно при выполнении предпосылок о распределении парных разностей; он не подменяет оценку размера эффекта. При расширении до нескольких рекомендаций одного пользователя повторная выборка выполняется на уровне пользователя, а не отдельных текстов. Остальные сравнения объявляются дополнительными; подтверждающие выводы требуют учёта множественности проверок.

Воспроизводимый пакет должен включать версию данных, временные границы, параметры рекомендателя, выбранные пары, профили и основания, промпты, точную версию LLM, настройки генерации, исходные ответы, инструкцию и разметку. Сгенерированные объяснения с разметкой составят собственный подготовленный набор данных. Пользовательское доверие, полезность и причинная верность внутреннему механизму рекомендателя данным протоколом не измеряются.

Представленный метод и план оценки задают проверяемую постановку для последующего исследования. До выполнения пилота нельзя утверждать повышение качества, переносимость на второй домен или превосходство над опубликованными системами.

\clearpage
\balance
\begin{thebibliography}{12}
\bibitem{survey} Zhang Y., Chen X. Explainable Recommendation: A Survey and New Perspectives // Foundations and Trends in Information Retrieval. 2020. Vol.~14, no.~1. P.~1--101. DOI: \url{https://doi.org/10.1561/1500000066}.
\bibitem{pepler} Li L., Zhang Y., Chen L. Personalized Prompt Learning for Explainable Recommendation // ACM Transactions on Information Systems. 2023. Vol.~41, no.~4. Art.~103. DOI: \url{https://doi.org/10.1145/3580488}.
\bibitem{upr} Ramos J., Rahmani H.~A., Wang X., Fu X., Lipani A. Transparent and Scrutable Recommendations Using Natural Language User Profiles // Proceedings of ACL. 2024. P.~13971--13984. DOI: \url{https://doi.org/10.18653/v1/2024.acl-long.753}.
\bibitem{xrec} Ma Q., Ren X., Huang C. XRec: Large Language Models for Explainable Recommendation // Findings of EMNLP. 2024. P.~391--402. DOI: \url{https://doi.org/10.18653/v1/2024.findings-emnlp.22}.
\bibitem{sosedka} Соседка А.~В., Ибрагимов П.~И. Количественный и качественный анализ аудитории Telegram в разрезе рекомендаций с использованием больших языковых моделей // Искусственный интеллект в промышленных, коммерческих, медицинских и финансовых приложениях: сборник статей семинара кафедры «Инженерной кибернетики». Вып.~2. М.: НИТУ «МИСИС», 2024. С.~116--120. URL: \url{https://sadekov.su/Articles/kik_2024_1.pdf}.
\bibitem{beresnev} Береснев Д.~В. Анализ подходов к использованию предобученных моделей в разработке корпоративных чат-ботов // Искусственный интеллект в промышленных, коммерческих, медицинских и финансовых приложениях: сборник статей семинара кафедры «Инженерной кибернетики». Вып.~1. М.: НИТУ «МИСИС», 2023. С.~68--77. URL: \url{https://sadekov.su/Articles/kik_2023.pdf}.
\bibitem{kozharinov} Марчукова А.~В., Кожаринов А.~С. Об одном подходе к прогнозированию развития организаций на основе эволюционного моделирования и методов принятия решений // Достижения учёных XXI века: материалы 5-й международной научно-практической конференции, 20 июля 2010~г. Тамбов: ТАМБОВПРИНТ, 2010. С.~47--56. Библиографические сведения: «Наука МИСиС 2010», электронное приложение, с.~19. URL: \url{https://misis.ru/files/-/474ab05797f7a203f177c0cec3bac04b/Nauka_MISIS_2010.pdf}.
\bibitem{kion} MTS Kion Implicit Contextualised Sequential Dataset for Movie Recommendation. Официальный репозиторий. URL: \url{https://github.com/irsafilo/KION_DATASET} (дата обращения: 07.11.2026).
\bibitem{amazon} Amazon Reviews 2023. McAuley Lab. URL: \url{https://amazon-reviews-2023.github.io/} (дата обращения: 07.11.2026).
\bibitem{als} Hu Y., Koren Y., Volinsky C. Collaborative Filtering for Implicit Feedback Datasets // Proceedings of ICDM. 2008. P.~263--272. DOI: \url{https://doi.org/10.1109/ICDM.2008.22}.
\bibitem{cohen} Cohen J. A Coefficient of Agreement for Nominal Scales // Educational and Psychological Measurement. 1960. Vol.~20, no.~1. P.~37--46. DOI: \url{https://doi.org/10.1177/001316446002000104}.
\bibitem{wilcoxon} Wilcoxon F. Individual Comparisons by Ranking Methods // Biometrics Bulletin. 1945. Vol.~1, no.~6. P.~80--83. DOI: \url{https://doi.org/10.2307/3001968}.
\end{thebibliography}
\end{document}
